% Separate refinement of the proved pointed witness; no new Laurent-family claim.
\begin{corollary}[Proved: unpointed formal-disc isomorphisms at the same precision]
\label{D:unpointed-Eisenstein-witness}
In Proposition~\ref{D:local-Eisenstein-linear-atlas}, keep the same
coefficient charts and the same integer $N$. In its isomorphism
witness omit the equation $b_{00}=0$, retaining the invertibility
of $b_{10}$ and the equations $g(\beta)=0\pmod{(t^N,f)}$.
For every $k$-algebra $R$, the modified witness has an $R$-point
over $(f,g)$ if and only if there is a continuous
$R[[t]]$-algebra isomorphism $B_g\simeq B_f$, without any
requirement on the specified augmentations to $R$.
The bounds of $7d+6$ coordinates and degree at most $e^2+1$
remain valid. This is a statement about the formal-disc models
already constructed, not a presentation theorem for all families
of finite etale $R((t))$-algebras.
\end{corollary}
\begin{proof}
Use the polynomial representative
\[
       \beta=\sum_{i<e,\,j<N}b_{ij}t^j\alpha^i\in B_f=R[[\alpha]]
\]
from (D.EA.7). The witness equations give $g(\beta)=t^N h$ for
some $h\in B_f$. Evaluation at $\alpha=0$, hence at $t=0$, gives
\[
                              b_{00}^{\,e}=0.           \tag{D.UE.1}
\]
The constant coefficient of $\beta$ as an $\alpha$-series is
$b_{00}$, and its linear coefficient is $b_{10}$, because
$t\in\alpha^eR[[\alpha]]$ and $e\ge2$.

We recall the elementary formal-coordinate fact over arbitrary
coefficient rings. A series with nilpotent constant coefficient
and invertible linear coefficient is a continuous automorphism
of a formal disc. To verify it, let $I$ be the nilpotent principal
ideal generated by its constant coefficient. The simple root
$0$ of its reduction modulo $I$ lifts, by the finite nilpotent
implicit-function recursion, to a root $a\in I$. Translation by
$a$ is a continuous formal-disc automorphism; its substitution
is well defined because $a$ is nilpotent. After that translation
the series has zero constant coefficient and invertible linear
coefficient, and the usual coefficient-by-coefficient inverse
applies. This proves the assertion, including continuity of the
inverse. Applied to (D.UE.1), it identifies $B_f$ with $R[[\beta]]$.

One must still justify that the Newton estimate uses the correct
parameter. Write $g(X)=X^e+tU_g(t,X)$. Its constant
$U_g(0,0)$ is invertible. Modulo $t$ the equation $g(\beta)=t^Nh$
says $\beta^e=0$, so $U_g(t,\beta)$ is a unit in $B_f$.
Since $N\ge2$, we obtain
\[
       \beta^e=t\bigl(-U_g(t,\beta)+t^{N-1}h\bigr)
                    =t\cdot\text{unit}.               \tag{D.UE.2}
\]
The $t$- and $\beta$-adic topologies are consequently cofinal.
The coefficient-weight conditions for $g$, now applied using
(D.UE.2), give
\[
                 g'(\beta)=\beta^d\cdot\text{unit},
       \qquad g(\beta)\in\beta^{eN}R[[\beta]],
       \qquad eN>2d.                                  \tag{D.UE.3}
\]
The all-coefficient-ring Newton argument in the proposition
therefore supplies an exact root $\beta_\infty$ with
\[
           \beta_\infty-\beta\in
                    \beta^{eN-d}R[[\beta]].            \tag{D.UE.4}
\]
Here $eN-d>d\ge1$, so $\beta_\infty$ is a pointed invertible
change of coordinate relative to $\beta$. Composing with the
already established, possibly unpointed, coordinate $\beta$
shows that $X\mapsto\beta_\infty$ defines a continuous
$R[[t]]$-algebra isomorphism $B_g\to B_f$.

Conversely, let such an isomorphism be given and let $\beta$ be
the image of the source coordinate. The source identity
$t=X^e\cdot\text{unit}$ implies that the constant coefficient
of $\beta$ is nilpotent. The inverse isomorphism is also
continuous and given by substitution of its source-coordinate
image. The chain rule for these mutually inverse substitutions
shows that the linear coefficient of $\beta$ is a unit: the
derivative of the inverse evaluated at the nilpotent constant
is defined and multiplies that coefficient to $1$.
Expand $\beta$ in the free $R[[t]]$-basis
$1,\alpha,\ldots,\alpha^{e-1}$ and truncate its coefficient
series modulo $t^N$. This retains its constant and linear
coefficients, gives the invertible $b_{10}$, and preserves
$g(\beta)=0$ modulo $t^N$. It is the required finite witness.
No additional variables or higher-degree equations have been
introduced.
\end{proof}
